% ============================================================
%  MODULE X — Advanced Topics & Next Steps
% ============================================================

\chapter{Advanced Topics \& Next Steps}

\begin{notebox}
This module introduces areas to explore as you develop mastery of the foundations.
Each topic below is a deep discipline in its own right. Treat these as
signposts for your learning roadmap, not immediate requirements.
\end{notebox}

% ────────────────────────────────────────────────────────────
\section{Options — The Basics}

An \textbf{option} is a contract giving the holder the \textbf{right, but not the obligation},
to buy or sell an asset at a specified \textbf{strike price} by or on the \textbf{expiration date}.

\begin{center}
\begin{tabular}{>{\bfseries}p{2.5cm} >{\bfseries}p{2.5cm} p{7.5cm}}
\toprule
Type & Position & Rights and Risk \\
\midrule
Call & Long & Right to \emph{buy} at strike. Profit if price rises above strike + premium. \\
Call & Short & Obligation to \emph{sell} at strike. Limited profit (premium), unlimited risk. \\
Put & Long & Right to \emph{sell} at strike. Profit if price falls below strike + premium. \\
Put & Short & Obligation to \emph{buy} at strike. Limited profit (premium), large downside. \\
\bottomrule
\end{tabular}
\end{center}

\subsection{The Greeks}
Options price sensitivity measures:

\begin{center}
\begin{tabular}{>{\bfseries}p{2cm} p{10.5cm}}
\toprule
Greek & What It Measures \\
\midrule
Delta ($\Delta$) & Change in option price per \$1 change in underlying. Ranges 0--1 (calls), 0 to $-$1 (puts). \\
Gamma ($\Gamma$) & Rate of change of Delta. High near expiration. \\
Theta ($\Theta$) & Time decay. Options lose value as expiration approaches. \\
Vega ($\mathcal{V}$) & Sensitivity to changes in implied volatility. \\
Rho ($\rho$) & Sensitivity to interest rate changes. \\
\bottomrule
\end{tabular}
\end{center}

\subsection{Basic Strategies}
\begin{itemize}
  \item \textbf{Covered Call} — own stock + sell call = generate income, cap upside
  \item \textbf{Protective Put} — own stock + buy put = insure downside
  \item \textbf{Straddle} — buy call + buy put at same strike = profit from big move either way
  \item \textbf{Vertical Spread} — buy one option + sell another at different strike = defined risk/reward
\end{itemize}

\begin{warningbox}
Options involve leverage and time decay. Short options have unlimited or very large risk.
Understand all mechanics before trading with real capital.
\end{warningbox}

% ────────────────────────────────────────────────────────────
\section{Futures Contracts}

A \textbf{futures contract} is a \textbf{standardised agreement} to buy or sell an asset
at a predetermined price on a specific future date.

Uses:
\begin{itemize}
  \item \textbf{Hedging} — lock in a price (e.g., airline hedging jet fuel)
  \item \textbf{Speculation} — directional bets with leverage
  \item \textbf{Arbitrage} — exploit price discrepancies between spot and futures
\end{itemize}

Key futures markets: \textbf{equity index} (S\&P 500 /ES, NASDAQ /NQ),
\textbf{bond} (US Treasury /ZB, /ZN), \textbf{energy} (Crude /CL, Natural Gas /NG),
\textbf{metals} (Gold /GC, Silver /SI), \textbf{currencies} (/6E, /6J).

\subsection{Contango and Backwardation}
\begin{itemize}
  \item \textbf{Contango}: Futures price $>$ spot price. Normal for storage-cost commodities.
        Rolling long positions forward incurs a cost.
  \item \textbf{Backwardation}: Futures price $<$ spot price. Indicates supply tightness.
        Rolling forward generates a positive roll yield.
\end{itemize}

% ────────────────────────────────────────────────────────────
% Term Cards / Glossary
% ────────────────────────────────────────────────────────────
\section{Trading Term Cards (Finance Jargon)}

\subsection{Assets and Cash Markets}

\begin{keyterm}{Equities}
Stocks / shares that represent ownership in a company. If the company grows, the equity
value can rise; if it deteriorates, it can fall.
\\\textbf{Example:} Buying 1 share of a tech company makes you a partial owner of its
future cash flows (via dividends and/or resale value).
\end{keyterm}

\begin{keyterm}{Fixed Income}
Securities where returns are driven primarily by interest payments and credit quality
(for example: government or corporate bonds). You earn coupons, and your total return
depends on yields moving after you buy.
\\\textbf{Example:} Buy a 5\% coupon bond; if market yields drop, the bond price can rise
even if the coupon stays the same.
\end{keyterm}

\begin{keyterm}{FX (Foreign Exchange)}
Trading currencies in pairs (e.g., \textbf{EUR/USD}). An FX quote is just the price of one
currency in terms of another.
\\\textbf{Example:} If EUR/USD moves from 1.1000 to 1.1050, EUR strengthened versus USD.
\end{keyterm}

\begin{keyterm}{Commodities}
Raw materials (oil, gold, wheat, copper, etc.) traded directly (spot) or via derivatives
(futures/options). Many commodities are priced in USD, so FX and macro matter a lot.
\\\textbf{Example:} Oil prices can jump on supply disruptions; gold often responds to real
rates and risk sentiment.
\end{keyterm}

\subsection{Derivatives}

\begin{keyterm}{Futures}
Exchange-traded, standardised contracts to buy or sell an underlying at a future date
for a predetermined price. Futures are marked-to-market (gains/losses flow daily).
\\\textbf{Example:} Hedge jet fuel costs for next quarter by locking in the fuel price
via futures.
\end{keyterm}

\begin{keyterm}{Options}
Contracts that give the holder the \textbf{right} (not obligation) to buy (call) or sell
(put) the underlying at a \textbf{strike} price on or before expiration.
\\\textbf{Example:} A call option benefits if the stock rises above strike; your downside
is limited to the premium you paid.
\end{keyterm}

\begin{keyterm}{Perp Swaps (Perpetual Futures Swaps)}
Derivative contracts with no fixed expiration date (``perpetual''). Traders typically
use them like futures, but the contract uses \textbf{funding} to keep the perp price close
to the relevant spot/mark price.
\\\textbf{Example:} If BTC perp trades too high versus BTC spot, funding tends to push the perp
back down by incentivizing more short demand.
\end{keyterm}

\begin{keyterm}{Funding Rate in Perp Swaps}
A periodic payment exchanged between longs and shorts. The sign/magnitude is designed to
reduce persistent price gaps (basis) between perp and spot.
\\\textbf{Example:} If perp price trades above spot (positive premium), longs typically pay
shorts (funding is positive for longs), discouraging over-longing.
\end{keyterm}

\begin{keyterm}{Forwards}
OTC (custom) contracts agreeing to buy/sell an asset at a future date for a fixed price.
Forwards are not standardized like exchange futures and therefore typically carry more
\textbf{counterparty risk}.
\\\textbf{Example:} A producer and airline might agree on fixed fuel pricing for next year
via an OTC forward.
\end{keyterm}

\subsection{Position Direction and Portfolio Risk}

\begin{keyterm}{Long}
Holding a position that benefits when the underlying price rises.
\\\textbf{Example:} Buying BTC at \$60{,}000 and selling at \$63{,}000 profits if your entry is lower
than your exit.
\end{keyterm}

\begin{keyterm}{Short}
Holding a position that benefits when the underlying price falls (via borrowing/selling
or via derivatives like futures/perps).
\\\textbf{Example:} If you short at \$100 and it falls to \$95, you profit (before fees/interest).
\end{keyterm}

\begin{keyterm}{Flat}
No active exposure: your net position is zero (or close enough that price moves do not
meaningfully affect your P\&L).
\\\textbf{Example:} After closing a long and a short, you are flat and stop taking directional risk.
\end{keyterm}

\begin{keyterm}{Exposure}
How sensitive your portfolio is to a specific risk factor (directional price moves, rates, FX,
volatility, etc.). Think of exposure as ``how much you win/lose if the variable moves.''
\\\textbf{Example:} Long 1 BTC means your mark-to-market moves roughly with BTC price by about 1:1
(ignoring funding/fees).
\end{keyterm}

\begin{keyterm}{Margin}
Collateral you post to support leveraged positions. Margin determines how much loss you
can absorb before a forced close.
\\\textbf{Example:} With \$1{,}000 initial margin, you might control \$10{,}000 notional (10x),
but you only have \$1{,}000 of buffer.
\end{keyterm}

\begin{keyterm}{Leverage}
Controlling a larger notional position relative to your posted equity/collateral.
Leverage amplifies both gains and losses.
\\\textbf{Example:} 10x leverage means a 1\% adverse move can cost about 10\% of your margin
(ignoring fees and mark-price mechanics).
\end{keyterm}

\begin{keyterm}{Liquidation}
The forced closing of a leveraged position when collateral/equity falls below the
maintenance requirement (or the exchange's liquidation threshold).
\\\textbf{Example:} BTC drops quickly; your margin cushion is depleted and the exchange closes you
at the prevailing mark/liquidation price.
\end{keyterm}

\begin{keyterm}{PnL (Profit and Loss)}
The amount you gain or lose from a trade. It can be realized (closed positions) or
unrealized (mark-to-market on open positions).
\\\textbf{Example:} Buy at \$100 and sell at \$110: realized P\&L is +\$10 per unit.
\end{keyterm}

\begin{keyterm}{Sharpe Ratio}
A performance metric that rewards higher returns while penalizing volatility:
\[
S = \frac{R_p - R_f}{\sigma_p}
\]
where $R_p$ = portfolio return, $R_f$ = risk-free rate, and $\sigma_p$ = return volatility.
\\\textbf{Example:} Two strategies both return 10\%; if one is much smoother (lower volatility),
it tends to have a higher Sharpe.
\end{keyterm}

\begin{keyterm}{Counterparty Risk}
Risk that the other side in a contract fails to perform (e.g., defaults on payments).
Exchange-cleared derivatives often reduce this risk compared to OTC contracts.
\\\textbf{Example:} An OTC forward depends on your counterparty paying/settling correctly.
\end{keyterm}

\subsection{Order Book, Execution, and Trading Costs}

\begin{keyterm}{Bid}
The highest price a buyer is willing to pay.
\\\textbf{Example:} If the bid is \$99.90, buyers are currently willing to buy at \$99.90.
\end{keyterm}

\begin{keyterm}{Ask}
The lowest price a seller is willing to accept.
\\\textbf{Example:} If the ask is \$100.10, sellers are currently offering to sell at \$100.10.
\end{keyterm}

\begin{keyterm}{Spread}
The difference between ask and bid. In liquid markets, the spread is your baseline trading
cost (you usually pay it when you enter/exit).
\\\textbf{Example:} Bid \$99.90, ask \$100.10: spread is \$0.20.
\end{keyterm}

\begin{keyterm}{Fill}
When an order is executed. Orders can fill partially (some quantity filled) or fully.
\\\textbf{Example:} Your \$10{,}000 limit buy might fill \$6{,}000 if remaining book liquidity
doesn't reach your price.
\end{keyterm}

\begin{keyterm}{Slippage}
The difference between the price you expected and the actual execution price, driven by
spread, fast market moves, and your order size/impact.
\\\textbf{Example:} You expect a fill near \$100 but the market moves and you end up filling at \$100.25.
\end{keyterm}

\begin{keyterm}{Market Order}
An order that executes immediately at the best available prices. You prioritize speed over
price certainty.
\\\textbf{Example:} In a thin order book, a market buy may walk up the book.
\end{keyterm}

\begin{keyterm}{Limit Order}
An order that executes only at your specified price or better. You prioritize price certainty over
guaranteed fill.
\\\textbf{Example:} A buy limit at \$99.50 won't buy if the ask never reaches \$99.50.
\end{keyterm}

\begin{keyterm}{Tick}
The smallest allowed change in price (minimum price increment).
\\\textbf{Example:} If tick size is 0.01 USDT on a venue, prices move in discrete steps of 0.01.
\end{keyterm}

\subsection{Market Makers, Inventory, and Crypto Basis}

\begin{keyterm}{Inventory}
The net positions market makers accumulate while providing two-sided liquidity.
Carrying inventory creates risk if the market moves against the position.
\\\textbf{Example:} If a market maker has sold more than bought recently, they may be net short and
can lose if price rises.
\end{keyterm}

\begin{keyterm}{Skew}
An asymmetry in pricing/quoting (often referring to options implied-volatility skew, or more generally
that market makers price different strikes differently).
\\\textbf{Example:} Downside puts can have higher implied volatility, creating a skew across strikes.
\end{keyterm}

\begin{keyterm}{Spread Capture}
How market makers try to profit: repeatedly buying at the bid and selling at the ask.
If execution is mostly from passive flow, this can be a steady return stream.
\\\textbf{Example:} Earn small profits on many trades while keeping your quotes tight.
\end{keyterm}

\begin{keyterm}{Adverse Selection}
The risk that trades hitting your quotes are informed (you get ``picked off'') rather than random.
It can turn spread capture into losses.
\\\textbf{Example:} You quote tightly, but sophisticated traders know price is about to move and
your filled trades become systematically loss-making.
\end{keyterm}

\begin{keyterm}{Basis}
The price gap between two related instruments (commonly: futures/perps vs spot).
In crypto, basis is tightly connected to funding.
\\\textbf{Example:} If BTC perp trades \$200 above BTC spot, positive basis creates incentive for funding
to move the perp back toward spot.
\end{keyterm}

\begin{keyterm}{Open Interest}
The total number of outstanding derivative contracts (not yet closed). It reflects how much positioning
is currently ``hanging out'' in the system.
\\\textbf{Example:} Open interest rising alongside price can indicate growing leverage; falling open interest
can indicate deleveraging.
\end{keyterm}

% ────────────────────────────────────────────────────────────
\section{Quantitative and Algorithmic Trading}

\subsection{Systematic vs. Discretionary}
\begin{itemize}
  \item \textbf{Discretionary} — human makes all trading decisions. Can incorporate
        qualitative/contextual reasoning.
  \item \textbf{Systematic} — rules-based; computer executes all trades. Removes emotion.
        Requires robust backtesting and forward testing.
\end{itemize}

\subsection{Strategy Development Process}
\begin{enumerate}[label=\textbf{\arabic*.}]
  \item \textbf{Hypothesis} — define the edge you believe exists and why it should persist.
  \item \textbf{Data collection} — quality historical data is non-negotiable.
  \item \textbf{Backtest} — simulate strategy on historical data. Watch for overfitting.
  \item \textbf{Walk-forward test} — test on data not used in development (out-of-sample).
  \item \textbf{Paper trading} — simulate live with real-time data, no real money.
  \item \textbf{Live trading (small size)} — deploy with minimal capital first.
  \item \textbf{Scale} — increase size only after demonstrable live performance.
\end{enumerate}

\begin{warningbox}
\textbf{Overfitting} is the cardinal sin of backtesting: a strategy that perfectly fits
historical data but fails on new data. The more parameters you optimise, the more you
risk fitting noise. Always hold out out-of-sample data for validation.
\end{warningbox}

\subsection{Common Quant Strategies}
\begin{itemize}
  \item \textbf{Momentum} — buy winners, short losers over 3--12 month lookbacks
  \item \textbf{Mean reversion} — bet on extreme deviations reverting to the mean
  \item \textbf{Statistical arbitrage (stat-arb)} — pairs or basket mean reversion
  \item \textbf{Trend following} — CTA style; long trending assets across all classes
  \item \textbf{Factor investing} — exposure to value, quality, size, low-volatility factors
\end{itemize}

% ────────────────────────────────────────────────────────────
\section{Behavioural Finance}

\begin{keyterm}{Behavioural Finance}
The study of how cognitive biases and emotional responses cause investors to make
irrational decisions, leading to systematic market mispricings.
\end{keyterm}

\begin{center}
\begin{tabular}{>{\bfseries}p{4cm} p{9cm}}
\toprule
Bias & Description \\
\midrule
Loss Aversion & Losses feel roughly $2\times$ worse than equivalent gains feel good \\
Overconfidence & Overestimating one's predictive accuracy and skill \\
Anchoring & Over-relying on the first piece of information seen (e.g., purchase price) \\
Recency Bias & Overweighting recent events; extrapolating the current trend \\
Herd Behaviour & Following the crowd; contributing to bubbles and crashes \\
Disposition Effect & Holding losers too long, selling winners too early \\
Sunk Cost Fallacy & Holding a bad investment to ``get back to even'' \\
Narrative Fallacy & Constructing plausible stories to explain random price moves \\
\bottomrule
\end{tabular}
\end{center}

% ────────────────────────────────────────────────────────────
\section{Recommended Learning Roadmap}

\subsection{Beginner (Months 1--3)}
\begin{itemize}[noitemsep]
  \item Complete this entire notes document
  \item Read: \textit{The Intelligent Investor} — Benjamin Graham
  \item Read: \textit{A Random Walk Down Wall Street} — Burton Malkiel
  \item Open a \textbf{paper trading account} (Interactive Brokers, Thinkorswim)
  \item Study and annotate 100 historical charts
\end{itemize}

\subsection{Intermediate (Months 4--9)}
\begin{itemize}[noitemsep]
  \item Read: \textit{Market Wizards} — Jack Schwager
  \item Read: \textit{Trading in the Zone} — Mark Douglas (psychology)
  \item Read: \textit{Technical Analysis of the Financial Markets} — John Murphy
  \item Develop and paper trade a \emph{single, well-defined} strategy for 3 months minimum
  \item Begin tracking trades in a \textbf{trading journal}
\end{itemize}

\subsection{Advanced (Month 10+)}
\begin{itemize}[noitemsep]
  \item Trade live, minimum size; focus on process, not P\&L
  \item Read: \textit{Reminiscences of a Stock Operator} — Edwin Lefèvre
  \item Study options via tastytrade or Options Education (CBOE)
  \item Explore macro with: \textit{The Alchemy of Finance} — George Soros
  \item Consider Python for backtesting: \texttt{pandas}, \texttt{vectorbt}, \texttt{backtrader}
\end{itemize}

% ────────────────────────────────────────────────────────────
\section{Essential Tools and Resources}

\begin{center}
\begin{tabular}{>{\bfseries}p{3.5cm} p{3cm} p{6.5cm}}
\toprule
Tool & Category & Use Case \\
\midrule
TradingView & Charting & Best charting platform; scripting via Pine Script \\
Interactive Brokers & Broker & Best for serious traders; global access, low fees \\
Bloomberg Terminal & Data/News & Institutional standard (expensive); use free alternatives \\
Yahoo Finance & Data & Free fundamental data and price history \\
EDGAR (SEC) & Filings & Official source for US company filings \\
Finviz & Screener & Stock screening by technical and fundamental filters \\
Quandl / FRED & Economic Data & Macro and economic time series data \\
Python / R & Quant & Backtesting, data analysis, strategy automation \\
\bottomrule
\end{tabular}
\end{center}

% ────────────────────────────────────────────────────────────
\section{Final Principles to Internalise}

\begin{enumerate}[label=\textbf{\Roman*.}]
  \item \textbf{Capital preservation first}. You cannot profit if you blow up.
  \item \textbf{Process over outcome}. A losing trade within your rules is fine.
        A winning trade outside your rules is dangerous.
  \item \textbf{The market is always right}. Your opinion is irrelevant; price is fact.
  \item \textbf{Cut losers, let winners run}. This is the opposite of what emotions dictate.
  \item \textbf{There is no holy grail}. Any edge erodes over time; keep learning.
  \item \textbf{Consistency beats brilliance}. A 60\% win rate with 2:1 RRR,
        repeated 1{,}000 times, compounds to large wealth.
  \item \textbf{Know your edge}. If you cannot articulate why your strategy works,
        you do not have one.
  \item \textbf{Keep a trading journal}. The data in your own trade history is the
        most valuable dataset you will ever have.
\end{enumerate}

\vspace{2em}
\begin{tcolorbox}[colback=lightblue, colframe=darkblue]
\centering
\textbf{``The goal of a successful trader is to make the best trades.\\
Money is secondary.''} — Alexander Elder
\end{tcolorbox}
